\documentclass[aspectratio=169]{beamer}
% \usetheme{dmpisa}
% \titlebackground*{assets/background}

\usetheme{Copenhagen}
\setbeamertemplate{navigation symbols}{}

% font settings
\usepackage{caladea}    % Cambria clone
\usepackage{carlito}    % Calibri clone
\renewcommand{\familydefault}{\sfdefault}
\usefonttheme[onlymath]{serif}

% This command is not strictly necessary but it allows you to add
% links in a nice color without having to set them each time
\newcommand{\hrefcol}[2]{\textcolor{cyan}{\href{#1}{#2}}}

% custom packages
\usepackage[italian]{babel}
\usepackage[T1]{fontenc}
\usepackage[utf8]{inputenc}

\title{Basi cristallo}
\subtitle{Un'applicazione alla corrispondenza  di Edelman--Greene}
% \course{Bachelor, master, Ph.D., and Research}
\author{\href{mailto:m.musumeci5@studenti.unipi.it}{Matteo Musumeci} \and Prof.~Michele D'Adderio}
% \IDnumber{0.57721$\ldots$}
\date{25 settembre 2026}

\begin{document}

\maketitle

% \section{Introduction}

\begin{frame}{Parole ridotte}
    \begin{columns}
        \column{.5\textwidth}
        \column{.5\textwidth}
        \begin{itemize}
            \item Posso solo scambiare due pedine vicine
            \item Devo usare il minimo numero di mosse
                (ovvero non posso scambiare le stesse due pedine più di una volta)
        \end{itemize}

    \end{columns}
\end{frame}

\begin{frame}{Parole ridotte}
    Ad ogni passo, annoto la posizione in cui è avvenuto lo scambio.

\end{frame}


\begin{frame}
    \begin{block}{Il problema}
        Fissata una permutazione,
        quante sono le sue parole ridotte?
    \end{block}

\end{frame}
\end{document}
